\section{Sandbox Guide and Development Workflow}

The \path{FSA_model_dev} repository is designed for iterative model development. This section outlines the workflow for developing and validating model changes.

\subsection{Environment Setup}

Development should be conducted within the \path{comfyenv} Conda environment, which contains the necessary JAX, Diffrax, and numerical libraries.
\begin{lstlisting}[language=bash]
conda activate comfyenv
export PYTHONPATH=.
\end{lstlisting}

\subsection{Running Simulations}

Forward simulations can be executed using the provided example script:
\begin{lstlisting}[language=bash]
python examples/run_fsa_simulation.py
\end{lstlisting}
This script simulates a 7-day training block followed by a recovery phase, generating a visualization of the fitness (\(B\)), fatigue (\(F\)), and autonomic (\(A\)) trajectories.

\subsection{Validation Suite}

Before a model is considered ready for the main SMC${}^2$ repository, it must pass two levels of validation:
\begin{enumerate}
    \item \textbf{Physics Check:} \path{tests/test_fsa_physics.py} verifies that the SDE integration is stable and that the state remains within physiological bounds.
    \item \textbf{Artifact Compatibility:} \path{tests/test_artifacts.py} ensures that the \path{EstimationModel} and \path{ControlSpec} can be instantiated and that their JAX-compiled functions are error-free.
\end{enumerate}

\subsection{Modifying the Model}

Researchers can modify the underlying physiology by editing \path{_dynamics.py}. Changes to the observation model or priors should be reflected in \path{estimation.py}. Once validated using the local test suite, the entire \path{models/fsa_high_res/} directory can be safely migrated to the production \path{python-smc2-filtering-control} repository.
